\chapter{What would count as success?}\label{ch:evidence}

\section{Four statements with different burdens of proof}
We can now state the core result without shrinking the ambition. For a declared complete physical transition and a fixed potential, its finite EBU is the potential before minus the potential after. With sufficient differentiability it is also the negative gradient integral along a declared path between those states. Consecutive actions under the same field telescope. These are mathematical statements under assumptions, not promises that a particular society has already been regulated.

The next statement is about an implementation. Software can check source limits, calculate a Gaussian quote exactly for rational inputs, preserve event provenance and reject inconsistent records. Passing those tests establishes conformance to the implemented contract. It does not show that the chosen field is a faithful measure of planetary function or that people will respond well to the incentives.

A third statement is scientific evidence. A registered study may compare EBU-guided action with controls under demand, disturbances, reserves and loss. It should publish every defined outcome, including failure and non-evaluable cases, under a frozen protocol. A demonstration that arithmetic closes is not evidence that a living system remains viable; it is a prerequisite for interpreting such evidence.

The final statement is institutional. To claim that an EBU economy improves long-term life support, one must examine who can participate, which harms cross its boundary, whether essential needs remain available and what happens outside the studied world. Those are not reasons to abandon the project. They identify the work required to make its promise real.

\section{What a long-run test must ask}
A genuine long-run study needs repeated physical opportunities or demands, replenishment or regeneration where those processes exist, disturbances and reserve constraints. It must specify before execution what counts as persistent restoring tendency or sustained function, and distinguish physical possibility from actor affordability. It should compare policies on the same external conditions without quietly giving one policy a better menu, a different arrival stream or a hidden optimizer.

At least three falsifiers are easy to name. The model may be physically incomplete: a positive EBU action might merely move harm outside the boundary. The policy may be ineffective: even with a faithful field and feasible choices, actors may not preserve service and reserves. The institution may be unfair: it might deny necessary action to those without historical capacity. A useful study is designed to expose these failures, not to make them impossible by definition.

\section{The special role of exactness}
The Gaussian accounting identities can be checked exactly under rational inputs. That is valuable because a failed equality then points to a model, implementation or record problem rather than an adjustable numerical tolerance. Yet exact arithmetic is not exact knowledge of the world. Measurements have error, reference values can be contested, and a formally complete state can still omit a relevant carrier. The right attitude is to be strict about the mathematics and candid about the empirical boundary.

\section{A future the account is meant to protect}
The planet will continue receiving sunlight and changing through cycles whether or not an accounting system exists. What is not guaranteed is that human action will preserve the living and built conditions on which future generations depend. A sale can be successful today while drawing down water, fertility, resilience or access. A repair can matter enormously even when no immediate sale rewards it.

EBU begins with the idea that physical balance should be visible in the decisions that shape those conditions. The field and finite action law make a precise form of that idea possible. They do not require the world to stand still at a perfect point. Homeostasis means sustaining function through change, including the ability to recover and adapt. That is why the ultimate measure of the project is not how many positive receipts it can issue, but whether a properly governed system helps people meet needs without consuming the basis of tomorrow's life.

This book has carried the reader from monetary incentives through physical laws, homeostasis, Gaussian geometry, exact finite accounting and the design of an inspectable institution. The route is continuous, but its destinations remain distinct: a theorem is proved, a mechanism is specified, an experiment is run, and a society decides whether the result is worth trusting. EBU's ambition deserves all four.
